\documentclass[11pt,letterpaper]{article}
\usepackage[T1]{fontenc}
\usepackage{lmodern,microtype}
\usepackage[margin=1in]{geometry}
\usepackage{amsmath,amssymb,booktabs,graphicx}
\usepackage[numbers,sort&compress]{natbib}
\usepackage{xcolor}
\usepackage[colorlinks=true,allcolors=blue!40!black]{hyperref}
\hypersetup{pdftitle={ALFM-BEM: Outcome-conditioned experience memory for frozen predictors},pdfauthor={David Ahmann}}
\setlength{\emergencystretch}{2em}
\title{\textbf{ALFM-BEM}\\[3pt]\large Outcome-conditioned experience memory for frozen predictors}
\author{David Ahmann\\ \small Independent researcher, Toronto, Canada}
\date{Open research preprint \quad$\cdot$\quad November 18, 2026}
\begin{document}
\maketitle
\begin{abstract}
A model operator with a limited review budget must choose which predictions to
accept and which to refer. Bidirectional Experience Memory (BEM) stores outcomes
alongside embeddings and retrieves evidence of both success and failure while
the underlying predictor stays fixed. We evaluate the memory against its
simplest controls. In a five-seed synthetic task, BEM obtains failure and
success F1 of 0.588 and 0.716. An outcome vote over ordinary cosine neighbors
obtains 0.744 and 0.734. A failure-only rule also reproduces the decisions of
the BEM-based filter in a simulated claims stream. We further evaluate review
allocation on held-out handwritten digits, with observed classifier errors,
identical review budgets and a change in the decision policy.
Across 12 paired replays, review of 720 out of 7,200 future case presentations catches 178 of 454 errors using the BEM risk score, 291 using classifier confidence, and 379 using a calibrated combination of confidence and a local class vote. A calibrated policy combining BEM risk, success, familiarity and confidence catches 290.

We release the implementation, per-seed results and executable protocols.
\end{abstract}

\section{The decision problem}
An operator can freeze a predictor yet change what happens to its predictions.
For example, if similar cases have repeatedly failed, the operator may refer
the next case for review. A failure among past accepted cases is an observation;
the outcome a withheld prediction would have produced can remain unknown.
Consequently, a fair evaluation must specify which decisions produce feedback
and charge each method for the reviews it requests.

ALFM-BEM (Adaptive Latent Feedback Model with Bidirectional Experience Memory)
was originally proposed as a wrapper for frozen foundation models \citep{ahmann}.
This study evaluates two narrower claims. First, how well do the memory scores
identify the outcomes of unseen cases relative to an outcome-labeled neighbor
vote? Second, when only a fixed number of reviews are available, which policy
sends the cases on which review will help most? An improvement on the first
question does not establish improvement on the second. In particular, a
classifier can predict failure accurately at the extremes yet rank the cases
near the review cutoff poorly.

\section{What the memory provides}
For a fixed representation $z$, the system stores normalized historical
embeddings, signed outcomes, timestamps and case identifiers. A repeated case
can update the corresponding entry. Separate read operations return supporting
successes and failures. A familiarity score describes the position of a case
relative to the stored embeddings. The public interface returns evidence and
case records, allowing a caller to inspect which observations influenced a
decision. Appendix~\ref{app:methodref} gives the score definitions and the
conditions under which an optional adapter update admits a parameter bound.

The full wrapper supports accepting a prediction, abstaining, escalating,
and requesting information. These actions incur different costs. Asking a
reviewer for a correct label can resolve an individual case immediately.
A claim about learning from that review requires evaluating later cases whose
labels were unavailable at decision time. The fixed-budget experiment isolates
the choice of cases for review; it does not simulate a conversation or train
a language model.

\section{Experiments and results}
\paragraph{A common review budget.}
The frozen predictor made 454 errors across 7,200 future case presentations. Each policy referred 720 cases. Figure~\ref{fig:review} counts the errors caught by a perfect reviewer. Each seed reuses a subset of the same finite corpus. The totals are paired replay counts; the same digit can appear in more than one replay.
\begin{figure}[htbp]\centering
\begin{tabular}{@{}lrl@{}}\toprule
Review ranking & Errors caught & Share of underlying errors\\\midrule
Arrival order & 41 & \textcolor{teal!70!black}{\rule{2.80mm}{1.25ex}}\hspace{0.45em}9.0\%\\
Nearest failure & 176 & \textcolor{teal!70!black}{\rule{12.02mm}{1.25ex}}\hspace{0.45em}38.8\%\\
BEM risk & 178 & \textcolor{teal!70!black}{\rule{12.15mm}{1.25ex}}\hspace{0.45em}39.2\%\\
Frozen confidence & 291 & \textcolor{teal!70!black}{\rule{19.87mm}{1.25ex}}\hspace{0.45em}64.1\%\\
Confidence and BEM risk & 291 & \textcolor{teal!70!black}{\rule{19.87mm}{1.25ex}}\hspace{0.45em}64.1\%\\
Joint BEM scores and confidence & 290 & \textcolor{teal!70!black}{\rule{19.80mm}{1.25ex}}\hspace{0.45em}63.9\%\\
Confidence and outcome vote & 319 & \textcolor{teal!70!black}{\rule{21.78mm}{1.25ex}}\hspace{0.45em}70.3\%\\
Confidence and class vote & 379 & \textcolor{teal!70!black}{\rule{25.88mm}{1.25ex}}\hspace{0.45em}83.5\%\\
\bottomrule\end{tabular}
\caption{Errors selected for review at identical capacity (five referrals per batch of fifty). The votes combine local evidence with classifier confidence using held-out calibration. Appendix~\ref{app:budgetref} gives the calibration protocol, uncertainty and cost-sensitive comparisons.}
\label{fig:review}\end{figure}

\paragraph{Outcome prediction.}
The binary classification metric used here scores the outcome of the incoming
query. It does not assess the usefulness of a retrieved precedent to a model
generating a new answer. In the original overlapping-cluster construction,
a symmetric success vote over the same neighbors used to vote on failures
has higher mean F1 for both outcomes. Representations, thresholds
and the generation of the cluster centers materially affect this comparison.
The executable controls in Appendix~\ref{app:controlref} preserve the
distinction between the historical setup and the controls.

\paragraph{Referral in a simulator.}
In the canonical 2,000-claim stream, 231 of 2,000 submitted claims would be
denied under always-submit. The memory-based filter submitted 1,775, of which
45 were denied, and referred 225. In the final 100 cases the comparable
figures are 16/100 and 1/84, with 16 cases referred. A failure-only
neighbor rule makes the same decisions on the five tested streams.
The simulated payer supplies counterfactual outcomes on withheld cases.
These counts characterize a synthetic referral policy and cannot be
used to estimate real claims performance or the cost of a human review.

\section{Relationship to prior work}
Selective classification explicitly trades coverage against risk
\citep{selective}. This distinction motivates our common budget and the
inclusion of a frozen-model confidence baseline. Memento learns to select
cases with a frozen language model \citep{memento}. ReasoningBank distills
strategies from successful and failed trajectories \citep{reasoningbank}.
Storing both outcome signs therefore has precedent. Our experiments isolate
the value of using their scores to allocate review.

Learning on the Job measures how a frozen agent uses outcome feedback and
corrections on later banking tasks \citep{job}. ReScopeBench tests the
maintenance and reuse of experience under silent environment changes
\citep{rescope}. Our controlled policy change similarly makes it possible
for old evidence to become stale. Our digit experiment uses a classifier and
a single explicit change. The agent studies evaluate additional capabilities,
including language-based rule induction and multi-step interaction.

\section{Limits and implications}
A memory score can be inspected without being calibrated as a probability.
A parameter norm bound controls a numerical update without establishing retention
of earlier capabilities. Accuracy after referral depends on who reviews the
case, how accurate that reviewer is and whether feedback returns. The experiment
assigns feedback and review outcomes explicitly. Its uncertainty estimates
describe variation across the specified repeated runs of one dataset.
Transfer to another model, task family or reviewer population requires new
evidence.

For deployment research, the next experiment is a prospective stream of
frozen-model decisions with actual task outcomes, delayed or missing feedback,
and recorded review time. The relevant comparison is whether memory
improves correct resolutions per incoming case at the same cost as a
confidence rule or an outcome vote. For the released implementation,
the present results establish reproducible reference points for that test.

\paragraph{Availability.}
The source, result records and build instructions accompany version 0.2.0 at
\url{https://github.com/davidahma/alfm-bem}. The digit images are distributed
through scikit-learn \citep{sklearn} from the optical recognition dataset
\citep{digits}. The other core controls generate synthetic records. An
optional text experiment evaluates keyword retrieval on 20 Newsgroups with
the published MiniLM sentence encoder \citep{newsgroups,minilm}.
The software is available under the MIT license. The paper and the author's
result tables are available under CC BY 4.0. AI-assisted coding and editing
tools were used to revise the implementation, develop comparative tests,
and prepare this release.

\bibliographystyle{unsrtnat}
\bibliography{references}
\appendix
\section{Memory scores, decision policy and adapter bounds}
\label{app:methodref}

\subsection{What enters the memory}

The memory accepts an embedding and an observed outcome $o\in[-1,1]$.
The full wrapper starts from a supplied backbone vector, applies a two-layer
ReLU projection and a chain of residual adapters, and stores the resulting
normalized vector. The projection starts with random weights. Its optional
offline training objective is a contrastive softmax loss on same-polarity
partners and opposite-polarity negatives. The direct memory interface accepts
a representation chosen by the caller. Neither entry point trains the backbone.

A record is keyed by case identifier, tenant and domain. The wrapper hashes
the supplied context to obtain the identifier. On a repeated key with $n$
previous observations, the memory increments the confirmation count, refreshes
the timestamp and replaces the outcome with $(n o_{\rm old}+o_{\rm new})/(n+1)$.
It also normalizes $n z_{\rm old}+z_{\rm new}$ after averaging. Confirmations
count submissions under a key and need not come from independent observers.
An outcome below $-0.3$ contributes failure evidence, and one above $0.3$
contributes success evidence. Every stored record contributes to familiarity.

\subsection{Evidence scores for one memory}

Let $s_i$ be the dot product of the normalized query and record $i$.
The implementation divides by the vector norm plus $10^{-10}$, so a zero
vector has zero similarity. Let $d_i=\exp(-\lambda a_i)$, where $a_i$ is
the integer age in days, and let $n_i$ be the confirmation count.
The severity $v_i$ is $5,4,3,2,1$ on the successive outcome intervals
$[-1,-0.8)$, $[-0.8,-0.5)$, $[-0.5,0)$, $[0,0.5)$ and $[0.5,1]$.

From failures with $s_i>\tau$, select up to 20 in decreasing similarity
order to form $F$. Write $[x]_a^b=\min(b,\max(a,x))$. The risk score is
\begin{align}
 w_i &= [\alpha\,s_i(v_i/5)d_i n_i]_0^{0.99},\\
 b_F &=
 \begin{cases}
  0.1\max_{i<j,\ i,j\in F} z_i^\top z_j & |F|\geq 2,\\
  0 & |F|<2,
 \end{cases}\\
 R &= \left[1-\prod_{i\in F}(1-w_i)+b_F\right]_0^1 .
\end{align}
For all successes $S$ with $s_i>\tau$, the success score is
\[
 U=\left[\frac{1}{|S|}\sum_{i\in S}s_i\,\frac{6-v_i}{5}\,d_i\right]_0^1 .
\]
An empty evidence set gives zero for its score. The two scores can both be
high. Shared neighbors, repeated confirmations and the pairwise term make
a probability interpretation of the noisy-OR unwarranted.

\subsection{Fixed-bandwidth familiarity}

For a memory with $N\geq 10$ records, the default coverage score uses bandwidth
$h=0.3$. Define
\[
 K_h(z,x)=\exp\!\left(-\frac{(1-z^\top x)^2}{2h^2}\right),\qquad
 D_N(z)=\frac{1}{N}\sum_{i=1}^N K_h(z,z_i).
\]
The memory evaluates $D_N$ at up to 256 stored vectors selected at evenly
spaced record indices, including the selected records' own contributions.
If $q$ is the 90th percentile of those self-scores, coverage is
$C(z)=[D_N(z)/\max(q,10^{-12})]_0^1$. Memories with fewer than ten
records return zero coverage. Rescaling is per memory and clipping can tie
queries with different unscaled scores.

For each fixed query $z$ and fixed $h>0$, if the stored $Z_i$ are independent
draws from a common distribution $P$ on the unit sphere, the bounded summands
obey the strong law of large numbers. Thus
\[
 D_N(z)\ \longrightarrow\ \int K_h(z,x)\,dP(x)
 \qquad\text{almost surely}.
\]
The limit measures kernel-weighted probability mass. For example, if
$P$ is a point mass at $u$, a query orthogonal to $u$ receives the positive
limit $\exp[-1/(2h^2)]$ despite its neighborhood having zero mass.
Identifying a pointwise density would require further assumptions, bandwidth
selection and an appropriate normalization. Deduplication, adaptive feedback
and pruning can also violate the sampling assumption. The result concerns the
unscaled kernel mean. It establishes neither convergence of the capped-sample
calibration quantile nor discrimination between arbitrary distributions.

\subsection{Aggregation, indexing and actions}

The manager reads the global memory, the requested domain memory if present,
and the requested tenant memory if present. It takes the maximum risk, maximum
success and arithmetic mean coverage across these memories, including an empty
global memory. The caller supplies trusted scope identifiers. This partitioning
rule assumes the application has already checked access to those scopes.

With a cached $N$-by-$d$ embedding matrix, a dense similarity pass takes
$O(Nd)$ arithmetic. Coverage also averages $N$ kernel values. Failure retrieval
can sort as many as $N$ candidates and computes the pairwise term on at most
20 vectors. Recalibrating after a write requires up to $\min(N,256)$ extra
similarity passes. The optional scikit-learn cosine neighbor index uses brute
force for unfiltered top-$k$ retrieval. Outcome-filtered retrieval and the
three decision scores use the dense matrix. There is no logarithmic query-time
guarantee for the released implementation.

Table~\ref{tab:release-policy} gives the first matching rule for the default
consensus engine. The user supplies deterministic heuristic checks. At fixed
risk, success and heuristic results, raising coverage can take a case out of
the Query region once. It cannot return the case to Query at higher coverage.
The numeric confidence attached to an action is a policy output.
\begin{table}[htbp]
\centering
\begin{tabular}{ll}
\toprule
Condition in priority order & Action\\
\midrule
Critical rule violated or failed to evaluate & Escalate\\
$C<0.2$ & Query\\
$R>0.85$ & Escalate\\
$R>0.6$ & Abstain\\
Noncritical rule violated & Abstain\\
$U>0.7$ & Trust\\
$C<0.5$ & Query for validation\\
Otherwise & Trust at low confidence\\
\bottomrule
\end{tabular}
\caption{Action thresholds and precedence. Equalities follow the next applicable rule.}
\label{tab:release-policy}
\end{table}

\subsection{What the adapter update bounds}

Let $\theta$ concatenate the trainable parameters of one residual adapter,
including the affine LayerNorm parameters and biases. For an embedding $z$,
the adapter normalizes $z+W_2\operatorname{ReLU}(W_1\operatorname{LN}(z)
+b_1)+b_2$. With $s$ the dot product of that output with the normalized replay
input, the example loss is $1-s$ for $o>0$ and $\max(0,s+0.5)$ otherwise.
The batch objective adds $\lambda_{\rm reg}\sum_j\|\theta^{(j)}\|_2$,
where each parameter tensor is flattened separately. The adapter clips
the gradient of the entire objective jointly, then applies plain SGD with
zero momentum and zero optimizer weight decay. The code reports the mean
example loss before regularization.

Let $B=\{v:\|v\|_2\leq r\}$ and $P_B$ be its Euclidean projection.
Initialization projects all parameters jointly before recording the baseline
or copying the exponential moving average. With a fixed radius, a finite
clipped gradient $\|g_t\|_2\leq c$ and step size $\eta_t\geq 0$, the update is
\[
 \theta_0=P_B(\theta_{\rm init}),\qquad
 \theta_{t+1}=P_B(\theta_t-\eta_t g_t).
\]
Since $P_B(\theta_t)=\theta_t$ and projection onto a closed convex ball
is nonexpansive,
\[
 \|\theta_{t+1}-\theta_t\|_2
 \leq\|(\theta_t-\eta_t g_t)-\theta_t\|_2\leq\eta_t c,\qquad
 \sum_{t=0}^{T-1}\|\theta_{t+1}-\theta_t\|_2\leq c\sum_{t=0}^{T-1}\eta_t .
\]
These bounds hold in exact arithmetic. The measured path length
aggregates the post-projection steps and excludes initialization. Endpoint
displacement is also bounded by $2r$. A raw weight import starts a new
measurement interval at its projected state. Loading an older adapter checkpoint
likewise establishes a fresh baseline for subsequent updates.

The moving average starts in $B$ and remains there by convexity for
$0\leq\beta\leq1$. Direct adapter calls use live weights by default and the
wrapper's inference chain uses the moving average. An output-displacement
bound at a fixed input would additionally require a Lipschitz bound in
\emph{parameters} along the relevant path. The parameter bound alone
establishes no decision agreement, task accuracy or retention guarantee.

\section{Fixed-budget review of handwritten digits}
\label{app:budgetref}

\paragraph{Cases and intervention.}
The task is to transcribe one handwritten digit. We use all 1,797 labeled
$8\times8$ images in scikit-learn's bundled \texttt{load\_digits} dataset.
The pixel arrays are distinct. A review consumes one slot, reveals the recorded
digit $y$, and replaces the proposed digit $a$ with $y$. An accepted case
retains $a$. We assume the reviewer returns the recorded digit without delay
or error. The feedback record keeps the original proposal and its corrected
label, so a reviewed mistake remains a failure of the original predictor.
We count both the slots consumed and the loss remaining after review. No
time or monetary value is assigned to a slot.

\paragraph{Frozen predictor and representation.}
A NumPy permutation with seed 4701 selects 397 training rows without reference
to the labels. We divide each pixel by 16 and fit multinomial logistic
regression with an intercept, $L_2$ penalty, $C=1$, the L-BFGS solver,
tolerance $10^{-8}$ and at most 500 iterations. The proposed digit is the
highest-probability class. For retrieval, we subtract the training-set mean
pixel vector and normalize each resulting vector to unit length. Every method
uses the same fixed proposals, class probabilities and representation. Training
rows serve only to fit the predictor and representation. The classifier made
92 errors on its 1,400 held-out images. All subsequent trials reuse this
model and pool.

\paragraph{Who sees which labels.}
For each seed from 0 through 11, a new label-independent permutation divides
the pool into four disjoint parts. The first 150 cases supply the same fully
reviewed initial memory to all methods. The next 150 are a fully labeled
calibration set, used only to set mixture weights and neighborhood sizes.
The following 500 cases arrive in ten batches of 50. Each method ranks a
complete batch, commits to five reviews, and receives only those five labels
before the next batch. Earlier arrivals break ties after scores are rounded
to 12 decimal places. Each adaptive history thus contains the same 150
initial labels and 50 additional labels selected by that method.

The final 600 cases arrive in twelve more batches. At this boundary we freeze
each history and evaluate two review policies on the same cases. Each policy
again receives five slots per batch. The evaluator uses the labels to count
losses and corrections. No future label updates a score, a weight or a memory.
The two future evaluations are counterfactual alternatives, each with 60
reviews. A replay allocates 410 post-training labels per policy (150 initial,
150 calibration, 50 stream and 60 future). The frozen classifier can operate
without the historical review labels. Equal access here isolates the choice
of cases to review at a given capacity.

\paragraph{Changing the cost of a mistake.}
During the stream every error costs one. We first retain this policy on the
future batches, then use the cost $|a-y|$ on those same batches. For example,
reading a 1 as a 9 now costs eight units, and reading an 8 as a 9 costs one.
The cost schedules are specified in advance and both can use the original
calibration labels. This models a change from class accuracy to accuracy of a
numerical entry. The distance cost is a chosen utility function. The observed
labels and classifier errors come from the dataset. A review changes the final
answer in the current case. Any benefit from acquiring feedback is measured
separately by scoring the same future batches with just the initial memory.

\paragraph{Comparators and calibration.}
Arrival order takes the first five cases in a randomly ordered batch.
Nearest failure ranks by maximum cosine similarity to a reviewed error.
The frozen rule uses $1-\max_c p(c\mid x)$ under unit loss. Under the
changed policy it uses $\sum_c p(c\mid x)|a-c|$, retaining the original
proposal $a$. Thus the classifier alone can respond to the cost change.

Two local votes use the same cosine neighbors and paid reviews. The outcome
vote averages the realized losses of the neighbors' own proposals. The
class-label vote estimates the digit distribution from their corrected labels
and applies the loss of the \emph{current} proposal to each of those labels.
For example, a correctly predicted neighboring 2 contributes zero past
error to the outcome vote, but supports review when the current proposal
is 7. A neighbor of cosine similarity $s$ has weight
$1/\max(1-s,10^{-8})$ before normalization. Each vote can combine its
estimate with the frozen score. For each objective separately, we choose
$k\in\{3,7,15\}$ and local weight $w\in\{0,0.25,0.5,0.75,1\}$ to minimize
mean squared error against realized loss on the 150 calibration cases,
using only the initial memory as neighbors. Ties favor smaller $k$, then
smaller $w$. These settings stay fixed for the rest of the replay.

\paragraph{Deployed BEM score and matched blends.}
We insert each reviewed embedding through
\texttt{BidirectionalExperienceMemory.add\_experience} and rank the batch by
its public \texttt{risk\_signal}. We use the package defaults for cosine
threshold (0.5) and risk sensitivity (0.8), with time decay disabled in
this undated replay. The risk score uses similar stored failures. Successes
are recorded but do not contribute to this read operation. The batch quota
is imposed by the experiment. We evaluate this risk component with a fixed
representation and without adapter updates.

For sign-only memory, success receives $+1$ and error receives $-1$.
For the changed cost, graded BEM re-encodes a reviewed error of distance $d$
as $-(0.35+0.65d/9)$. All nine possible distances pass the package's failure
threshold. Its severity buckets map distances 1--2, 3--6 and 7--9 to
severities 3, 4 and 5, respectively. We replay the exact same reviewed
records into the sign and graded memories. The encoding is computed entirely
from the paid historical labels.

We also give BEM the same opportunity to use the frozen score as either vote.
For both encodings we minimize the same calibration loss over the same five
weights in $wR+(1-w)F$, with $R$ the BEM score and $F$ the frozen score.
Under distance loss we first multiply $R$ by the mean distance of errors
in the initial memory, common to both encodings, to express the blend in
digit-distance units. The fallback scale for an error-free initial memory is
one. The sign and graded blends share the review history collected by the
unit-loss blend. Raw BEM collects its own history during the stream.
Comparisons between those two histories measure the complete review policy.

\paragraph{Joint use of failure, success and familiarity.}
A separate policy makes its own five stream reviews per batch. It calls the
same memory's \texttt{risk\_signal}, \texttt{success\_signal} and
\texttt{coverage\_signal} on each case when the calibrated joint weight is
positive. Write their outputs as $R,S,C$. Default coverage is the
package's calibrated cosine-kernel density score. We rank by
\[
  w h \{\alpha R+\beta(1-S)+\gamma(1-C)\}+(1-w)F .
\]
Here $h=1$ under unit loss and equals the initial mean distance per error
under distance loss. We use sign feedback for the unit objective and
re-encode the same policy history with graded outcomes for distance.
Calibration considers four ordered triples $(\alpha,\beta,\gamma)$, starting with
equal thirds, then $(1/2,1/4,1/4)$ and its two permutations emphasizing
lack of success or lack of coverage. Each component contributes when $w>0$.
We search the same five outer weights $w$ as the other blends, using all
20 combinations on the same calibration cases and initial memory. Ties
favor smaller $w$ and then the first triple in the order given. At $w=0$
the classifier alone ranks reviews, so the inner weights have no effect.
This is one specified way of combining the three package signals.

\begin{table}[htbp]
\centering
\caption{Paid reviews and their downstream outcomes. The stream columns cover
600 reviews and 412 underlying errors in 6,000 presentations. Each future
policy has 720 reviews, 454 underlying errors and 2,074 available
digit-distance units in 7,200 presentations. Counts pool 12 paired replays.}
\label{tab:digit-budget}
\small
\begin{tabular}{@{}lrrrrr@{}}
\toprule
 & Stream & \multicolumn{2}{c}{Future, unit cost} &
   \multicolumn{2}{c}{Future, distance cost} \\
\cmidrule(lr){3-4}\cmidrule(l){5-6}
Ranking & Caught & Caught & Left & Saved & Left \\
\midrule
Arrival order             &  33 &  41 & 413 &  184 & 1,890 \\
Nearest failure           & 143 & 176 & 278 &  944 & 1,130 \\
BEM risk, sign            & 140 & 178 & 276 &  959 & 1,115 \\
BEM risk, graded          & 140 & 178 & 276 &  975 & 1,099 \\
Frozen predictor          & 259 & 291 & 163 & 1,349 & 725 \\
BEM sign + frozen         & 251 & 291 & 163 & 1,454 & 620 \\
BEM graded + frozen       & 251 & 291 & 163 & 1,455 & 619 \\
Joint BEM + frozen        & 258 & 290 & 164 & 1,353 & 721 \\
Outcome vote + frozen     & 257 & 319 & 135 & 1,531 & 543 \\
Class-label vote + frozen & 303 & 379 &  75 & 1,693 & 381 \\
\bottomrule
\end{tabular}
\end{table}

\paragraph{Performance on future decisions.}
Table~\ref{tab:digit-budget} counts corrected errors under unit cost and
distance avoided under the changed cost. The class-label vote caught
379/454 errors (83.5\%), compared with 291/454 (64.1\%) for frozen confidence
and 178/454 (39.2\%) for raw BEM. Of the 6,480 automatically accepted
predictions, 75, 163 and 276, respectively, were wrong. With 720 perfect
reviews included, the corresponding end-to-end accuracies were 98.96\%,
97.74\% and 96.17\%. These totals count case presentations since each image
can recur in other seeds. The raw BEM advantage over nearest failure was
0.17 errors caught per seed (95\% interval $[-1.17,1.58]$).

Updating the class-label vote with its 50 paid stream labels changed 155 of
its 720 future review slots relative to its initial-memory version and
caught 41 more future errors, or 3.42 per seed $[1.83,5.00]$.
Raw BEM changed 200 slots and caught 25 more, or 2.08 per seed
$[0.83,3.33]$. In the matched BEM blend, 160 slots changed but the total
caught fell by three, a difference of $-0.25$ per seed $[-1.67,0.92]$.
Acquiring labels can change a review decision without improving its outcome.

\paragraph{What blending and grading changed.}
Under unit loss, calibration assigned positive weight to BEM in ten seeds
($w=0.25$ in seven and $w=0.5$ in three). The matched blend and frozen
confidence caught the same total, with a paired difference of 0.00 errors
per seed $[-2.25,2.08]$. The class-label and outcome votes caught 7.33
$[5.08,9.83]$ and 2.33 $[0.42,4.67]$ more errors per seed than the matched
BEM blend.

Under the changed cost, re-encoding raw BEM displaced 60 of 720 reviews
and saved 16 additional distance units, or 1.33 per seed $[-1.67,4.50]$.
The signed and graded BEM blends each used positive BEM weight in nine
seeds. The graded blend saved 106 more distance units than the frozen
expected-loss rule, or 8.83 per seed $[2.50,15.42]$, and one more unit
than the signed blend, or 0.08 per seed $[-2.83,3.50]$. The class-label
vote saved 19.83 more units per seed than the graded blend $[13.58,27.42]$.
The data support a gain over the frozen cost-sensitive rule for this blend
in these replays. The incremental contribution of graded encoding is
unresolved.

\paragraph{How often the joint score affected review.}
Calibration selected the classifier alone in nine of twelve unit-loss
replays. The joint weight was 0.25 twice and 0.5 once. Under distance loss
the joint weight was zero in eleven seeds and 0.75 in one. In every active
blend the inner weights were $(0.5,0.25,0.25)$ on failure risk, lack of
success and lack of coverage. At the same capacity, the policy caught
290 errors and avoided 1,353 distance units. Its paired differences from
the frozen rule were $-0.08$ errors per seed $[-0.50,0.25]$ and $+0.33$
distance units $[0.00,1.00]$. The risk-only graded blend saved 8.50 more
distance units per seed $[1.58,15.33]$ than this joint rule. These data
give little support for adding the three signals in this particular
combination to classifier confidence.

\begin{table}[htbp]
\centering\small
\caption{Paired future results by replay seed. Under unit loss the entries
count corrected errors out of the available errors. The final columns give
extra distance units saved by the graded BEM blend against each comparator
under the changed cost. Every cell uses the same budget of 60 reviews per method
and seed.}
\label{tab:digit-paired}
\begin{tabular}{@{}rrrrrrrr@{}}
\toprule
 & & \multicolumn{4}{c}{Errors caught} & \multicolumn{2}{c}{Graded blend gain} \\
\cmidrule(lr){3-6}\cmidrule(l){7-8}
Seed & Errors & Frozen & Risk blend & Joint blend & Class vote & vs.\ frozen & vs.\ sign blend \\
\midrule
0 & 40 & 28 & 23 & 28 & 30 & 0 & 0 \\
1 & 35 & 22 & 27 & 22 & 31 & 10 & 4 \\
2 & 42 & 20 & 23 & 20 & 36 & 26 & -6 \\
3 & 38 & 27 & 26 & 27 & 33 & 4 & -9 \\
4 & 36 & 21 & 26 & 21 & 32 & 30 & 0 \\
5 & 44 & 28 & 28 & 28 & 28 & 0 & 0 \\
6 & 32 & 24 & 20 & 22 & 28 & -13 & -5 \\
7 & 35 & 21 & 25 & 21 & 31 & 8 & 0 \\
8 & 35 & 23 & 24 & 23 & 29 & 13 & 3 \\
9 & 40 & 26 & 18 & 26 & 35 & 18 & 15 \\
10 & 31 & 21 & 21 & 21 & 26 & 0 & 0 \\
11 & 46 & 30 & 30 & 31 & 40 & 10 & -1 \\
\bottomrule
\end{tabular}
\end{table}

\paragraph{Uncertainty, scope and reproduction.}
Table~\ref{tab:digit-paired} shows how the effect varies by replay.
Each interval is the 2.5th and 97.5th percentile of 10,000 bootstrap means
of paired differences, resampling the twelve whole seeds with replacement
(NumPy seed 15001). All methods share the same batches within a draw.
The intervals describe sensitivity to random partitions and arrival order
on this finite corpus, conditional on the one trained classifier. They give
no independent replication across populations, writers or trained models.
The arrivals have simulated order, the reviewer is perfect, and the changed
cost is announced. There is no observed production shift, human review cost,
or evidence about language-model responses in this experiment.

From the release root with its environment active, reproduce the per-seed
counts, calibrated weights, selection comparisons and paired intervals with
\begin{verbatim}
python -m experiments.budget_benchmark --check results/budget.json
\end{verbatim}
The check uses the packaged digit data, one numerical thread and the
repository's BEM implementation. It requires exact counts and review-order
hashes, and tolerates floating-point roundoff in the reported summaries.

\section{Controls for outcome prediction, familiarity and referral}
\label{app:controlref}

\subsection{A symmetric vote on the original synthetic task}

Each memory contains 500 failures and 500 successes in 64 dimensions. The
test set contains 100 fresh examples of each sign drawn around the same
centers. For each sign, 70 examples come from its own centers with Gaussian
coordinate noise of standard deviation 0.05. The other 30 come from the
opposite centers with noise of standard deviation 0.10. The training set has
the same proportions. All observed vectors are normalized after adding noise.

The original generator draws ten failure centers and five success centers,
then divides each of the two matrices by its Frobenius norm. The mean squared
length of a center is consequently about $1/10$ for failure centers and $1/5$
for success centers. The same coordinate noise obscures a typical failure
center more strongly. A second run normalizes each center to unit length
before adding noise, using the same random draws in the same order. This
changes both the relative signal strengths and the absolute cluster separation.

The neighbor vote takes the ten highest cosine similarities in the mixed
memory, keeping those above 0.5. Its failure and success scores are the
fractions of those neighbors with the corresponding outcome sign. Both
scores are zero for an empty neighborhood. We flag a sign when its score
exceeds 0.3. BEM uses the same similarity and decision thresholds.
The failure-only comparator flags a failure for any stored failure above
the similarity cutoff and always emits zero success evidence. Binary F1
scores the sign of the incoming case. A query can cross both thresholds.

\begin{center}
\small
\begin{tabular}{lrrrr}
\toprule
 & \multicolumn{2}{c}{Original centers} & \multicolumn{2}{c}{Unit centers} \\
Method & Failure F1 & Success F1 & Failure F1 & Success F1 \\
\midrule
Neighbor vote & $0.744\pm0.021$ & $0.734\pm0.016$ & 0.700 & 0.700 \\
Failure only  & $0.568\pm0.022$ & $0.000\pm0.000$ & 0.667 & 0.000 \\
BEM           & $0.588\pm0.033$ & $0.716\pm0.013$ & 0.667 & 0.667 \\
\bottomrule
\end{tabular}
\end{center}

Entries give the mean and population standard deviation over seeds 42, 123,
456, 789 and 1011. The unit-center columns have zero observed dispersion.
In that setting BEM flags every query on both signs. Each task has 100 true
positives and 100 false positives, giving F1 of $2/3$. On the original
geometry the vote has higher mean F1 on both signs. Ordinary cosine retrieval
supports a positive vote given access to the stored outcomes.

\paragraph{Coverage on fresh in-distribution queries.}
For each seed the in-distribution (ID) queries are 200 additional draws from
the core failure centers. One out-of-distribution (OOD) set contains 200
vectors near a single new direction. Another contains 200 normalized
isotropic Gaussian draws. AUROC is the probability that a random ID score
exceeds a random OOD score, with half credit for a tie. The table averages
the five seed-level AUROCs on the original geometry.

\begin{center}
\begin{tabular}{lrr}
\toprule
Familiarity score & Clustered OOD & Uniform OOD \\
\midrule
Maximum cosine, all experiences & 0.9911 & 0.9885 \\
Maximum cosine, failures only   & 0.9920 & 0.9917 \\
BEM coverage, bandwidth 0.3     & 0.8792 & 0.9552 \\
\bottomrule
\end{tabular}
\end{center}

With unit centers, every entry is 1.0000. The ID sampling rule favors cases
close to a known failure mode. These AUROCs measure familiarity for that
population. Downstream error detection would require observed errors.

\paragraph{A threshold can hide separation.}
In a separate 768-dimensional control, seed 42 generates ten unit failure
centers, adds coordinate noise of standard deviation 0.06 to failures, and
draws successes isotropically. Maximum cosine to a stored failure has AUROC
1 on 100 fresh failures and 100 fresh successes. A cutoff of 0.50 detects
0 failures, 0.30 detects 94, and both 0.25 and 0.20 detect all 100. All four
cutoffs flag zero successes. The raw vectors already separate the classes
under this generator.

\subsection{Which claims were referred?}

A simulated payer applies three deterministic rules involving procedure,
diagnosis, age and modifier. A fixed random embedding sums vectors for the
procedure, diagnosis, modifier and decade of age. The payer uses exact age,
so two cases in one age bucket can have different outcomes. Each policy
acts before seeing the current outcome. BEM refers when risk exceeds 0.6
using a neighbor cutoff of 0.85. The failure-only comparator refers when
any previously observed denial has cosine similarity above 0.85.

\begin{center}
\small
\begin{tabular}{llrrrr}
\toprule
Seed 42 & Policy & Submitted & Denied on submit & Referred & Payable referrals \\
\midrule
All 2,000 & Always submit & 2,000 & 231 & 0 & 0 \\
          & Either filter & 1,775 & 45  & 225 & 39 \\
Final 100 & Always submit & 100   & 16  & 0 & 0 \\
          & Either filter & 84    & 1   & 16 & 1 \\
\bottomrule
\end{tabular}
\end{center}

Under either filter the full-stream denial rate among submissions is
$45/1775=2.54\%$. The filter refers 186 claims that would be denied and
39 that would be paid. Over the final 100 incoming claims the comparison
is $16/100$ versus $1/84$, with 16 referrals. The simulator supplies the
counterfactual outcome on a referral for evaluation.

We replay each of the five 2,000-claim streams twice. In the first replay
each agent observes every outcome after acting. In the second it learns only
from its own submissions. BEM and the failure-only rule agree on all 10,000
decisions within each replay. Feedback availability still matters. At seed
789 the two policies each refer 220 cases under full feedback and 199 under
submission-only feedback. Payable referrals fall from 35 to 14, while 45
submitted claims are denied in each replay. Review cost, recovery after
referral and real payer performance require other observations.

\subsection{Querying an oracle}

The original query simulation draws 1,000 Gaussian points in two dimensions
at seed 42. Labels indicate the interior of the unit circle. Prediction
correctness is sampled from a known confidence function that rises from
0.5 on the boundary to 1.0 at radial distance 0.2 from it. The same predictions
underlie all rows. A referral either removes a case from the automatic
denominator or supplies its true label immediately.

\begin{center}
\small
\begin{tabular}{lrrrr}
\toprule
Rule & Referred & Errors referred & Automatic accuracy & With oracle, all cases \\
\midrule
Trust all & 0 & 0 & $934/1000$ & $934/1000$ \\
Confidence below 0.7 & 109 & 48 & $873/891$ & $982/1000$ \\
Lowest 150 confidences & 150 & 58 & $842/850$ & $992/1000$ \\
\bottomrule
\end{tabular}
\end{center}

The 150 queries correct 58 of the original 66 errors. At either referral
budget, deferring those same cases to an equally accurate reviewer gives
the same all-case result. The oracle answers only the referred case.
Measuring what a system learns from review would require an update and
observations on subsequent predictions.

\subsection{Text controls with a frozen encoder}

The optional CPU experiment uses frozen 384-dimensional MiniLM sentence
embeddings. On the authored fixture, memory holds seven examples per outcome
sign. Three more of each sign are held out, with five separately authored
OOD sentences. Raw BEM coverage has AUROC 1.000 when the ID queries are
the stored sentences themselves and 0.800 on the six held-out ID sentences.
The 64-dimensional projection fits five sentences of each sign and uses
two per sign for validation. Its corresponding coverage AUROCs at seed
42 are 1.000 and 0.533. The raw failure score flags none of the six held-out
cases and the projected score flags all six. Only 30 ID--OOD pairs contribute
to each held-out AUROC.

\paragraph{A corpus with two different kinds of label.}
The 20 Newsgroups by-date corpus supplies observed topic labels. We use
space, medicine and electronics for the science domain, and baseball and
hockey for the sports domain. The historical failure proxy labels a science
document positive if its lowercased text contains \texttt{orbit} or
\texttt{circuit}. It measures substring retrieval. Topic discrimination
uses the newsgroup names. Measuring a predictor's failure on a user request
would require a separate outcome label.

We remove headers, footers and quoted text with the scikit-learn loader and
split each science outcome sign separately at seed 42. Of 4,954 documents,
255 keyword and 2,113 other science documents form the training pool. The
original balanced index contains all 255 keyword documents and 255 other
documents. A separate 64-dimensional projection is trained for 20 full-batch
updates on this balanced subset. The 593 science test documents contain 64
keyword cases and 529 others. All 1,993 sports documents are outside the
training domain. The stricter evaluation also removes blank bodies, bodies
seen anywhere in science training and repeated evaluation bodies, comparing
the stripped strings exactly. This leaves 64 keyword, 510 other science and
1,931 sports queries. Related posts with different text can still occur in
both training and evaluation.

\begin{center}
\small
\begin{tabular}{llrrr}
\toprule
Index & Representation & Keyword queries & Other science & Keyword AUROC \\
\midrule
Balanced & Raw       & 0.731 & 0.398 & 0.774 \\
Balanced & Projected & 0.950 & 0.764 & 0.682 \\
Full     & Raw       & 0.406 & 0.097 & 0.787 \\
Full     & Projected & 0.756 & 0.388 & 0.783 \\
\bottomrule
\end{tabular}
\end{center}

The middle columns give the mean fraction of keyword documents among five
cosine neighbors on unique unseen bodies. Projection raises this fraction
for both query classes. With the balanced index, ranking the keyword label
becomes less accurate even as retrieval on keyword queries alone rises.

\begin{center}
\small
\begin{tabular}{lrrr}
\toprule
Representation & Maximum cosine & Euclidean log KDE & Public BEM coverage \\
\midrule
Raw       & 0.946 & 0.857 & 0.897 \\
Projected & 0.735 & 0.692 & 0.691 \\
\bottomrule
\end{tabular}
\end{center}

These AUROCs compare the 510 unique non-keyword science bodies with 1,931
sports bodies, using the same balanced index throughout. The Euclidean
Gaussian KDE uses unit vectors and bandwidth 0.5. Public BEM coverage uses
its calibrated and clipped cosine-distance kernel at bandwidth 0.3.
On the unfiltered document split, raw maximum cosine scores 0.919 and
projected Euclidean log KDE scores 0.683. Comparing only those two numbers
changes the representation and the estimator together. The full result file
reports all six combinations for both evaluation populations.

\paragraph{Reproduction.}
From the repository root with the project environment active, run
\begin{verbatim}
python -m experiments.controls --output results/controls.json
python -m experiments.controls --check results/controls.json
\end{verbatim}
The second command reruns every synthetic control and compares all fields.
Counts must match exactly. Floating-point tolerance is $10^{-7}$ absolute
or $10^{-6}$ relative. These CPU controls use the project dependencies,
including NumPy and scikit-learn. The reference run used Python 3.13.7,
NumPy 2.3.5 and scikit-learn 1.7.2 on an Apple M4 Pro and took about five
seconds per command.

The optional text run obtains its corpus and weights from the
\href{https://scikit-learn.org/stable/modules/generated/sklearn.datasets.fetch_20newsgroups.html}{scikit-learn 20 Newsgroups loader}
and the
\href{https://huggingface.co/sentence-transformers/all-MiniLM-L6-v2/tree/c9745ed1d9f207416be6d2e6f8de32d1f16199bf}{pinned MiniLM model revision}.
The release contains aggregate measurements. To populate local caches and
reproduce the text results, run
\begin{verbatim}
python -m pip install -c constraints/research.txt -e '.[text]'
python -m experiments.real_embeddings --output results/text.json
\end{verbatim}
After the first run, add \texttt{--offline} to require cached inputs.
The reference run used Python 3.14.0, PyTorch 2.9.1, Sentence Transformers
5.1.2 and Transformers 4.57.3, taking about 24 seconds with cached inputs on
the same CPU. The result file records the model revision, numerical package
versions and a SHA-256 digest of the ordered topic labels and processed texts.

\end{document}
